% LIVE-READY STAGING ONLY, 28 August 2026.
% Release gate: do not integrate this prose until the arbitrary-field
% nonreduced exact-live statement/proof bytes receive exact-live verification.
% The rendered prose below deliberately contains no gate-status language.

\paragraph{From local flux to a global current.}
Proper tropical pushforward sends a localized positive closed current to a
positive closed Euclidean current supported on a one-dimensional polyhedral
set.  On the interior of an edge, positivity and the two closedness equations
force a constant nonnegative multiple of the integration current.  At a
vertex the endpoint terms satisfy the full balancing equation; no ray of the
active target star may be discarded.

Divide the edge flux of the given current by the calibrated flux of the
integration current.  The expansion-index comparison makes this ratio
intrinsic.  At a type-\(3\) point, the two opposite rays have the same ratio.
At a type-\(2\) point, write \(\operatorname{Rat}(C_x)\) for the function
field of the residual curve \(C_x\).  Appending a lift of
\(g\in\operatorname{Rat}(C_x)^\times\) turns the scalar component of
balancing into
\begin{equation}
        \sum_{p\in C_x} a_p m_p\operatorname{ord}_p(g)=0,
 \label{eq:p21p3-principal-balance}
\end{equation}
where \(m_p>0\) is the calibrated branch flux and \(a_p\) is the intrinsic
ratio.  Equivalently, define additive functionals on the divisor group by
\[
                        L_T([p])=a_pm_p,
                        \qquad L_Y([p])=m_p.
\]
Both are nonnegative on effective divisors, and
\eqref{eq:p21p3-principal-balance} says that \(L_T\) annihilates principal
divisors; the same balancing calculation for the integration current gives
this property for \(L_Y\).  Riemann--Roch over the actual residue field
forces each functional to be a nonnegative multiple of degree.  Comparing
the two multiples shows that all \(a_p\) are equal, without a rational point
on \(C_x\) and without algebraic closedness of \(K\).

The ratios propagate through all type-\(2\) and type-\(3\) points of a
connected rig-smooth curve.  The remaining type-\(1\) and type-\(4\) set has
tropical dimension zero.  By the tropical-dimension-zero test-form theorem,
the germ of every ordinary smooth \((1,1)\)-test form vanishes at each of
its points.  A current remainder already zero off that set therefore pairs
to zero with every test form.  This endpoint argument is distinct from the
rigid-point zero-atom assertion: the latter is used only at the locally
finite rigid conductor points when one descends through normalization.

This proves the connected classification in
\cref{prop:p21p3-gfc-smooth-rigidity}.  Passing through the finite
normalizations of the irreducible components gives
\eqref{eq:p21p3-curve-classification}.  Finally,
\cref{lem:p21p3-gfc-rigid-zero-atom} prevents the positive curvature
measures on different normalization sheets from cancelling at the
conductor.  That last step is why
\cref{prop:p21p3-curve-rigidity} proves Bedford--Taylor independence for
continuous psh pairs on the reduced curve itself, rather than only on its
normalization.  The nilreduction result
\cref{prop:p21p3-curve-rigidity-nonreduced} then shows that the ordinary
CLD calculus and the continuous-pair independence property are unchanged by
a nilpotent thickening, while the fundamental integration current retains
the generic component multiplicities.
